% Exact LaTeX recreation of the supplied Phase-1.pdf
% Compile with pdfLaTeX. Place member pictures in the same folder if required.

\documentclass[11pt,a4paper]{article}

\usepackage[a4paper,left=2.0cm,right=2.0cm,top=1.55cm,bottom=1.55cm]{geometry}
\usepackage[T1]{fontenc}
\usepackage{lmodern}
\usepackage{graphicx}
\usepackage{xcolor}
\usepackage{array}
\usepackage{tabularx}
\usepackage{ragged2e}
\usepackage{enumitem}
\usepackage{fancyhdr}
\usepackage{tikz}
\usepackage{setspace}

% -------------------------------------------------
% COLOURS
% -------------------------------------------------
\definecolor{TitleBlue}{RGB}{61,103,154}

% -------------------------------------------------
% GENERAL SETTINGS
% -------------------------------------------------
\setlength{\parindent}{0pt}
\setlength{\parskip}{0pt}
\renewcommand{\arraystretch}{1.0}

% The supplied document uses a clean sans-serif heading style
% and italic serif-style explanatory text.
\newcommand{\blueheading}[1]{%
    {\sffamily\color{TitleBlue}\fontsize{15}{18}\selectfont #1}%
}

\newcommand{\bluebigheading}[1]{%
    {\sffamily\bfseries\color{TitleBlue}\fontsize{16}{19}\selectfont #1}%
}

\newcommand{\instruction}[1]{%
    {\itshape\fontsize{10.5}{12.5}\selectfont #1}%
}

\newcommand{\answerbox}[2][3.0cm]{%
    \vspace{2pt}
    \noindent
    \fbox{%
        \begin{minipage}[t][#1][t]{\dimexpr\textwidth-2\fboxsep-2\fboxrule\relax}
            \vspace{1mm}
            \itshape #2
        \end{minipage}
    }
}

% -------------------------------------------------
% HEADER / FOOTER
% -------------------------------------------------
\pagestyle{fancy}
\fancyhf{}
\renewcommand{\headrulewidth}{0pt}
\renewcommand{\footrulewidth}{0pt}

\fancyhead[L]{%
    \sffamily\bfseries\itshape\fontsize{10.5}{12}\selectfont
     DSCPP-III-2026-T320
}
\fancyhead[R]{%
    \sffamily\bfseries\fontsize{10.5}{12}\selectfont
    PROJECT PROPOSAL
}
\fancyfoot[R]{%
    \itshape\fontsize{10}{12}\selectfont
    Page \thepage
}
\setlength{\headheight}{14pt}
\setlength{\headsep}{23pt}
\setlength{\footskip}{24pt}

% -------------------------------------------------
% PAGE 1
% -------------------------------------------------
\begin{document}

\vspace*{4mm}

\bluebigheading{PROJECT AND TEAM INFORMATION}

\vspace{13mm}

\blueheading{\textbf{Project Title}}

\vspace{1mm}

\answerbox[1.0cm]{ \textit{Rescue Food And Management System} }}

\vspace{13mm}

\blueheading{Student / Team Information}

\vspace{2mm}

% Team information table
% Compact 4-member layout so all members remain on Page 1.
\noindent
\begin{tabularx}{\textwidth}{|>{\RaggedRight\arraybackslash}p{0.69\textwidth}|>{\RaggedRight\arraybackslash}X|}
\hline
\begin{minipage}[t]{\linewidth}
\vspace{1.5mm}
\textit{Team Name:}\\[-1mm]
\textit{Team \#}
\vspace{1.5mm}
\end{minipage}
&
\begin{minipage}[t]{\linewidth}
\vspace{1.5mm}
\textit{
ResQbite
}
\vspace{1.5mm}
\end{minipage}
\\
\hline
\begin{minipage}[t][3.25cm][t]{\linewidth}
\vspace{1.5mm}
\textbf{Team member 1 (Team Lead)}\\[-1mm]
\textit{(Singh, name: Akshat student ID: 2510380179 2510380179@geu.ac.in, picture):}
\vspace{2.0cm}
\end{minipage}
&
\begin{minipage}[t][3.25cm][t]{\linewidth}
\vspace{1.5mm}
\textit{Singh,Akshat  -- 2510380179}\\
\textit{2510380179@geu.ac.in}

\vspace{2mm}
\includegraphics[width=3.0cm,height=1.5cm,keepaspectratio]{member1.png}
\end{minipage}
\\
\hline
\begin{minipage}[t][3.25cm][t]{\linewidth}
\vspace{1.5mm}
\textbf{Team member 2}\\[-1mm]
\textit{(Raturi, Hardik: 2510380179: 2510380179@geu.ac.in, picture):}
\vspace{2.0cm}
\end{minipage}
&
\begin{minipage}[t][3.25cm][t]{\linewidth}
\vspace{1.5mm}
\textit{Raturi,Hardik -- 2510380119}\\
\textit{2510380119@geu.ac.in}

\vspace{2mm}
\includegraphics[width=3.0cm,height=1.5cm,keepaspectratio]{member2.png}
\end{minipage}
\\
\hline
\begin{minipage}[t][3.25cm][t]{\linewidth}
\vspace{1.5mm}
\textbf{Team member 3}\\[-1mm]
\textit{(Gupta, Aadi: 2510388268: 2510388268@geu.ac.in, picture):}
\vspace{2.0cm}
\end{minipage}
&
\begin{minipage}[t][3.25cm][t]{\linewidth}
\vspace{1.5mm}
\textit{Gupta,Aadi -- 2510388268}\\
\textit{2510388268@geu.ac.in}

\vspace{2mm}
\includegraphics[width=3.0cm,height=1.5cm,keepaspectratio]{member3.png}
\end{minipage}
\\
\hline



\\

\end{tabularx}

\newpage

% -------------------------------------------------
% PAGE 2
% -------------------------------------------------

\bluebigheading{PROPOSAL DESCRIPTION (10 pts)}

\vspace{12mm}

\blueheading{Motivation (1 pt)}

\vspace{1mm}


\answerbox[6.0cm]{Foods that could have been eaten are being thrown away by restaurants, hotels, events, supermarkets, and homes every day, while a lot of people and organizations are still unable to get enough food. The main reason for this situation is the ineffective process of surplus food collecting, tracking, matching, and distributing. Food donorships are done manually, so that is why they are inefficient in tracking inventory, responding to requests quickly and correctly, and do not let people know about urgency requests and find necessary recipients for food. 

\vspace{1mm}FoodRescue is meant to solve the problem of food shortages by connecting food donators with people and organizations in need of food such as non-profit organizations, shelters, and other community service organizations. Donators will register available food with its specifications, and recipients will request food through the system. The system will provide priority for requests, food matching, transportation organization, and record keeping. 

\vspace{1mm}The aim of the system is to reduce food wasting.  }

\vspace{8mm}

\blueheading{State of the Art / Current solution (1 pt)}

\vspace{1mm}


\answerbox[3.5cm]{
At present, the surplus food is largely dealt with via phone calls, messaging applications, spreadsheets or personal communication between donors and NGOs. Though some organizations use digital systems, the relevant information may still be disparate and stored in different electronic formats.

\vspace{1mm}This makes tracking available food, its amount, and expiry dates very difficult. Consequently, food donations are often delayed, misallocated or lost entirely.

\vspace{1mm}FoodRescue streamlines the process by offering one single system for managing donors, food inventory, recipients, requests, food matching, and deliveries.
}

\vspace{7mm}

\blueheading{Project Goals and Milestones (2 pts)}

\vspace{1mm}


\answerbox[7cm]{
FoodRescue's primary objective is to minimize food waste by facilitating the sharing of excess food supplies between food donors and food recipients. Food inventory, requests, matching and deliveries will be handled by the system with the help of C++, DSA and SQLite. The main objectives are to sign up donors and recipients, maintain surplus food and stock, prioritize and match food requests, update food status and track deliveries, and effectively use DSA and OOP concepts.
### Milestones
\vspace{1mm}The project will start with the Requirement Analysis to determine its features and workflow. This will be followed by a System Design, where the System Architecture and Database Design will be developed. Next, the required SQLite tables will be implemented in **Database Setup**. The food, request, matching and delivery systems will then be developed in the Backend Development phase. Then, the process of **DSA Integration** will be done by using appropriate data structures and algorithms. The Frontend Integration will be used to link the frontend to the backend. Last but not least, there will be a Testing and Finalization stage that will test the system, fix bugs, and prepare the document.
\vspace{1mm}The ultimate objective is to create a **FoodRescue prototype suitable for use** that will enable the more efficient and structured redistribution of surplus food.
 }

\vspace{30mm}

\blueheading{Project Approach (3 pts)}

\vspace{1mm}

\instruction{}

\answerbox[9.1cm]{
The proposed solution is the FoodRescue Management System that will facilitate the collaboration between food donors and food beneficiaries like food banks, shelters, and NGOs in order to diminish food waste. The donors will register and give information on surplus food and recipients will submit food requests. The system will coordinate food available and prioritise food and allocate them based on need and food request.

\vspace{1mm}The application will be the object-oriented application and will have following components **Donors, Food Items, Recipients, Requests, and Deliveries** with the major part of the application to be developed in a backend language **C++**.

\vspace{1mm}A lightweight, easy-to-setup database will be used, SQLite. It will contain information about the donors, inventory of food and requests, and information about the food that is sent, and the C++ will perform operations with the database.

\vspace{1mm}These concepts will be used as needed: Priority Queues will be used for urgent requests, Normal Queues will be used for normal requests and deliveries, Dynamic Records will be managed with Linked List, and Undo or Cancel operations will be handled with stacks.

\vspace{1mm}The first version will be a console-based C++ Application developed in Visual Studio Code or another C++ IDE. The idea is to build a working prototype that showcases C++ programming, OOP, DSA and SQLite, and eventually provide a graphical or web interface if necessary.
}

\newpage

% -------------------------------------------------
% PAGE 3
% -------------------------------------------------

\blueheading{System Architecture (High Level Diagram)(2 pts)}

\vspace{1mm}

\answerbox[18.0cm]{
\includegraphics[width=18.0cm,height=100cm,keepaspectratio]{architecture.png}
}
\newpage

\blueheading{Project Outcome / Deliverables (1 pts)}

\vspace{1mm}

\answerbox[6cm]{
The FoodRescue Management System will offer a working solution for managing surplus food donations and linking them with the recipients (NGOs, shelter, food bank). The system will enable the users to enter the donors and recipients, add and manage available food, create food requests, prioritise the food requests and match the available food with the food needs of the recipients.

The project will be implemented with C++ language with the help of various data structures like linked list, queue, priority queue and stack using object-oriented principles. Donor, food, request and delivery records will be stored permanently in SQLite.

\vspace{1mm}Results include automated food allocation, updated inventory in real time, priority based requests, tracking of deliveries, expiry monitoring and reporting on food rescued, distributed, and wasted. The system will also be used for efficient application of DSA and OOP concepts in real-life applications. Final deliverables will include all source code, SQLite database schema, project documentation and working prototype.

}

\vspace{8mm}

\blueheading{Assumptions}

\vspace{3mm}


\answerbox[3.5cm]{

The project is based on the assumption that the donors give correct information regarding the type, quantity, condition and the expiry of the food. It also makes the assumption that there are legitimate food needs of those who are registered and that the food you are providing is safe to eat. Donors, recipients, and delivery volunteers will need to have contact and location details. The system assumes for simplicity that basic delivery routes and transport is available. Decisions about the safety of food will be made based on data entered in the system and simple validation rules - not by a live inspection or an advanced tracking system.
}

\vspace{12mm}

\blueheading{References}

\vspace{1mm}


\answerbox[4cm]{[1]:Do VIR: The Virtualization of Food Donation Distribution — An online system capable of connecting surplus food donors to recipients and increasing the effectiveness in distributing funds.
 \vspace{1mm}

[2]:A Web-Based Food Donation Distribution System for Effective Surplus Food Management — An online platform which provides surplus food donation management service by connecting the relevant surplus food to the respective recipients.

\vspace{1mm}
[3]:A Web Platform for Food Donation Distribution — A service created to reduce and minimize food waste by connecting food donors to various charity organizations supplying their surplus food donation.}

\end{document}
